\section{Discussion, Limitations, and Conclusion}\label{sec:conclusion}

\subsection{Discussion: routing as the unit of analysis}

The dominant framing in retrieval-augmented question answering, ``one method
wins,'' collapses several distinct failure modes into a single accuracy delta
and produces recipes that do not transfer across benchmarks. Our
bottleneck-aware view inverts that framing: \emph{routing} is the unit of
analysis, and each benchmark is characterized by which intervention helps,
which is neutral, and which actively hurts. Under this view, negative
controls are not embarrassments to be hidden but signal. Fixed
\hyre{} underperforming baseline on HousingQA and CaseHOLD is not
a failure of \hyre{}. It is positive evidence about which bottleneck each
benchmark exposes (metadata scope and statutory entailment on HousingQA,
answer-option conversion on CaseHOLD), and therefore a guide to the
intervention each benchmark actually needs.

Reading the calibration ladder this way also relocates the methodological
novelty of the \hyre{}/HyRE family. The two-call
architecture is a small efficiency variant of HyDE \citep{gao2023hyde}; what
is new is the \emph{snap-conditioned} passage prompt, which yields a
controlling-rule embedding rather than a question paraphrase, and the
placement of that embedding inside a routing portfolio rather than as a
universal default. The same HyDE-family intervention is positive on BarExam
(calibration \hyre{} $86.0$), neutral on LegalBench-SCALR ($76.0$
vs.\ baseline $74.0$), and negative on HousingQA ($51.5$ vs.\ $60.5$) and
CaseHOLD ($72.0$ vs.\ $73.0$). The accuracy signature is a property of the
\emph{bottleneck}, not the method.

The diagnostic protocol behind the \controller{} is dataset- and
model-agnostic by construction: any benchmark slice combined with any
language model and any retrieval stack produces a calibration trace whose
signals can be read off without retraining or annotated complexity
labels. This stands in contrast to learned routers
\citep{asai2024selfrag,jeong2024adaptive}, which require either
dataset-specific fine-tuning of reflection or critique tokens or a
supervised notion of question complexity. The protocol is therefore a
candidate \emph{measurement} layer rather than a competing recipe; we expect
it to compose with learned components in future systems rather than replace
them.

\subsection{Limitations}

Several limitations qualify the scope of these claims.
\begin{enumerate}
\item \textbf{Slice sizes.} The calibration tables use $N{=}200$ per
benchmark and the held-out tables $N{=}50$. The four legal benchmarks have
full sizes $1{,}195$ / $6{,}853$ / $3{,}600$ / $571$. Larger-$N$
evaluations are in progress on each benchmark and the numbers in this draft
should be read as diagnostic-grade evidence at fixed slice size, not final
benchmark rankings.
\item \textbf{Cross-model coverage.} The Llama~3.3~70B sanity sweep is
$N{=}50$ per benchmark and only partially transfers: the HousingQA
verifier and SCALR frontier routes hold up directionally, BarExam
\method{Snap-HyRE-v2} underperforms baseline ($72.0$ vs.\ $76.0$), and one
CaseHOLD diverse-HyRE cell did not produce a usable detail log on our run.
Multi-architecture generalization is therefore weakly, not strongly,
supported.
\item \textbf{Rule-based controller.} The \controller{} is rule and
evidence-summary based rather than learned. A learned router could plausibly
exploit calibration-trace features that the current rule set cannot read
off directly, and the gap between the rule-based controller and a learned
upper bound is itself an interesting open question.
\item \textbf{Held-out instability.} The held-out slice is rows $200$--$249$
($N{=}50$ per dataset). BarExam ($76.0$ selected vs.\ $90.0$ rewrite) and
SCALR ($80.0$ selected vs.\ $84.0$ frontier) are route-unstable on this
slice. The routing policy is informative but not finished; a rewrite-vs-HyRE
selector and a frontier-aware SCALR rule are the obvious next steps.
\item \textbf{Reference baselines.} We do not yet compare against full
\citet{asai2024selfrag} Self-RAG, \citet{yan2024crag} CRAG, or
\citet{wang2025specrag} Speculative RAG on legal benchmarks. Those
comparisons require either retraining or per-question critique
infrastructure and are deferred to future work.
\end{enumerate}

\subsection{Conclusion}

Legal RAG fails through bottleneck-typed mechanisms rather than through a
single missing capability. A cheap calibration trace, a fixed set of
signals collected on a $200$-row slice per benchmark, is enough to expose
which mechanism is active. A rule-based \controller{} that routes among
baseline retrieval, query rewrite, \hyre{}/HyRE,
state-filter retrieval, option grounding, a conservative verifier,
disagreement arbitration, and reject/escalate lifts macro accuracy by
roughly $6$\pp{} across the four legal benchmarks at a sub-$2$ LLM-call
budget. \hyre{} is one route among many: useful where the
bottleneck is query or legal-reasoning formulation, ineffective or harmful
where the bottleneck is metadata scope or answer-option conversion. The
contribution of this paper is therefore not a new universal recipe but a
discipline, \emph{name and measure the bottleneck before paying for the
intervention}, and the \framework{} controller that operationalizes that
discipline.
